\documentclass[11pt,a4paper]{article}
\usepackage[utf8]{inputenc}
\usepackage[T1]{fontenc}
\usepackage{lmodern}
\usepackage{amsmath,amssymb,amsfonts,amsthm}
\usepackage{geometry}
\geometry{margin=1in}
\usepackage{xcolor}
\usepackage{hyperref}
\usepackage{cleveref}
\usepackage{booktabs}
\usepackage{array}
\usepackage{cite}
\usepackage{microtype}
\usepackage{tcolorbox}

\hypersetup{
    colorlinks=true,
    linkcolor=blue,
    citecolor=red,
    urlcolor=cyan,
    pdftitle={Mathieu Moonshine Computed, the Double-Scaled Little String Bridge, and Dyons on K3 x T2 in Lean 4},
    pdfauthor={Xavier Callens and the SocrateAI Scientific Agora Collaboration}
}

\newtheorem{theorem}{Theorem}[section]
\newtheorem{lemma}[theorem]{Lemma}
\newtheorem{definition}[theorem]{Definition}
\newtheorem{proposition}[theorem]{Proposition}
\newtheorem{corollary}[theorem]{Corollary}
\newtheorem{conjecture}[theorem]{Conjecture}

\theoremstyle{remark}
\newtheorem{remark}[theorem]{Remark}

% ---------------------------------------------------------------------------
% Epistemic tier badges, following papers 7-9.
% ---------------------------------------------------------------------------
\newcommand{\tierA}{\colorbox{green!15}{\textsf{\scriptsize Tier A --- Lean kernel}}}
\newcommand{\tierL}{\colorbox{blue!10}{\textsf{\scriptsize Tier L --- literature}}}
\newcommand{\tierC}{\colorbox{orange!15}{\textsf{\scriptsize Tier C --- conjecture}}}
\newcommand{\lean}[1]{\texttt{#1}}
\newcommand{\Mtf}{M_{24}}

\newtcolorbox{scopebox}{%
  colback=gray!5, colframe=gray!45, boxrule=0.4pt, arc=1pt,
  left=6pt, right=6pt, top=3pt, bottom=3pt,
  fonttitle=\bfseries\small, title={Scope of Lean 4 verification}
}
\newtcolorbox{casebox}[1]{%
  colback=red!4, colframe=red!45!black, boxrule=0.4pt, arc=1pt,
  left=6pt, right=6pt, top=3pt, bottom=3pt,
  fonttitle=\bfseries\small, title={#1}
}

\title{\Huge \textbf{Mathieu Moonshine Computed,}\\[0.3em]
\LARGE \textbf{the Double-Scaled Little String Bridge,}\\[0.2em]
and Dyons on $K3\times T^2$ in Lean~4}

\author{\textbf{Xavier Callens}$^{1}$ \and \textbf{SocrateAI Scientific Agora Collaboration}$^{1}$ \\[0.5em]
$^{1}$\textit{SocrateAI Research Laboratory for Mathematical Physics \& Formal Verification}}

\date{September 2026}

% [layout] allow long Lean identifiers to break after underscores
\makeatletter\let\LMoldunderscore\_\renewcommand{\_}{\LMoldunderscore\allowbreak}\makeatother
\setlength{\emergencystretch}{3em}
% [/layout]
\begin{document}

\maketitle

\begin{abstract}
We report two Lean~4 libraries (toolchain \texttt{v4.33.1}, Mathlib) in which the mathematics
around the $K3$ elliptic genus is \emph{computed} from closed formulas rather than transcribed from
tables. \texttt{DualScaleMoonshine} (Stream~4) computes Mathieu moonshine from both sides: the
mock modular form $H^{(2)}=(-2E_2+48F_2)/\eta^3$ and its 21 twined series from their formulas, and
the traces of all 26 conjugacy classes of $\Mtf$ from the character table over $\mathbb{Z}[b_7]$,
$\mathbb{Z}[b_{15}]$, $\mathbb{Z}[b_{23}]$; the two meet at every class, and the first ten graded
pieces of the moonshine module decompose, from computed data, exactly as in Cheng--Duncan--Harvey's
printed table. The same library locates the ``24'' of the shadow in the polar/finite decomposition of
the elliptic genus, and reproduces the arithmetic of Harvey--Murthy--Nazaroglu's double-scaled little
string theories, including a two-sided proof of a divisibility they left unexplained.
\texttt{DualScaleDyons} (Stream~5) builds the quarter-BPS dyon partition function $1/\Phi_{10}$ on
$K3\times T^2$ from the computed elliptic genus, and shows through the orders stated that its
single-centred (``immortal'') part at $m=1,2,3$ is given by Hurwitz class numbers, and that its
$\Mtf$-twisted versions have virtual-character coefficients. The twining (``forger's'') test, applied
throughout, rejects an integer coincidence proposed earlier in this programme (the ``27720 lock'')
at all 25 non-identity classes, while $24=\chi(K3)$ and the Göttsche numbers pass it. Four
reading notes on the literature are recorded. Everything computed is Tier~A (kernel-checked; 101
and 44 theorems, 0 failing; axioms \texttt{propext}, \texttt{Classical.choice}, \texttt{Quot.sound}
at most); every physical reading is Tier~L or Tier~C. These are results about the mathematics of
$\mathcal{N}=4$ string compactifications on $K3\times T^2$, not about our universe: the one
experimental prediction of this programme has separately been found excluded by existing data.
\end{abstract}

\newpage
\tableofcontents
\newpage

%%%%%%%%%%%%%%%%%%%%%%%%%%%%%%%%%%%%%%%%%%%%%%%%%%%%%%%%%%%%%%%%%%%%%%%%%%%%%%%
\section{Introduction}
\label{sec:intro}
%%%%%%%%%%%%%%%%%%%%%%%%%%%%%%%%%%%%%%%%%%%%%%%%%%%%%%%%%%%%%%%%%%%%%%%%%%%%%%%

\subsection{Computed, not typed}
Earlier libraries of this programme checked Mathieu moonshine as one typed table against another:
the coefficients $A_1,\dots,A_9$ of Eguchi--Ooguri--Tachikawa~\cite{EOT2010} against the dimensions
of the irreducible representations of $\Mtf$. A transcription matching a transcription is weak
evidence --- and this project has already caught one mistyped $\Mtf$ table. The two libraries reported
here replace transcriptions by computations. The coefficients come from a closed formula; the twined
coefficients from the formula for $H_g$; the group-theoretic traces from the character table; the
dyon degeneracies from a Borcherds product whose only input is the $K3$ elliptic genus computed from
Jacobi theta functions. Where two independent computations meet, the coincidence is checked by the
Lean kernel, not copied. Where a printed table is used, it is used as a \emph{transcription control}
that a misreading would break.

\subsection{The forger's test}
What separated monstrous moonshine from numerology was that the coincidence survives \emph{twining}:
replacing a dimension by the trace of every group element. We apply the same test to this programme's
own integers (\Cref{sec:forger}). A relation proposed in paper~7~\cite{Callens2026Paper7}
(Theorem~6.1, the ``27720 lock'') fails it at every non-identity class; $24=\chi(K3)$ and the
numbers built from it pass it everywhere they are tested.

\subsection{Contents}
\Cref{sec:tiers} restates the tier system. \Cref{sec:moonshine} computes Mathieu moonshine from both
sides. \Cref{sec:forger} applies the forger's test. \Cref{sec:shadow} locates the ``24'' of the
shadow. \Cref{sec:hmn} reproduces Harvey--Murthy--Nazaroglu's double-scaled little strings.
\Cref{sec:dyons} builds the dyon partition function on $K3\times T^2$ and its single-centred part.
\Cref{sec:twisted} treats the $\Mtf$-twisted dyons. \Cref{sec:notes} records four reading notes on
the literature. \Cref{sec:context} places these results against the programme's experimental verdict.
\Cref{sec:scope} states what is and is not proved; \Cref{sec:methods} records methods;
\Cref{sec:open,sec:directions} give open problems and research directions; \Cref{app:repro} is a
reproducibility appendix.

%%%%%%%%%%%%%%%%%%%%%%%%%%%%%%%%%%%%%%%%%%%%%%%%%%%%%%%%%%%%%%%%%%%%%%%%%%%%%%%
\section{Epistemic tiers}
\label{sec:tiers}
%%%%%%%%%%%%%%%%%%%%%%%%%%%%%%%%%%%%%%%%%%%%%%%%%%%%%%%%%%%%%%%%%%%%%%%%%%%%%%%
As in papers~7--9, every claim carries one of three tags.
\begin{itemize}
  \item \tierA\ --- accepted by the Lean~4 kernel: \texttt{\#print axioms} reports at most
  \texttt{propext}, \texttt{Classical.choice}, \texttt{Quot.sound}; no \texttt{sorry}, no
  \texttt{admit}, no project-declared \texttt{axiom}, and no compiled evaluation
  (\texttt{native\_decide} is not used anywhere).
  \item \tierL\ --- established in the cited literature, quoted with a file and line range in
  \texttt{papers/foundations/*.txt}.
  \item \tierC\ --- a reading or conjecture introduced by this project. \textbf{The identification of
  a Lean object with a physical quantity is never Tier~A.}
\end{itemize}
Nearly all Tier~A statements below are finite identities between truncated power series with integer
(or Laurent-polynomial) coefficients, decided by kernel evaluation; each theorem states its truncation
order. ``Through $q^N$'' means exactly that and nothing beyond.

%%%%%%%%%%%%%%%%%%%%%%%%%%%%%%%%%%%%%%%%%%%%%%%%%%%%%%%%%%%%%%%%%%%%%%%%%%%%%%%
\section{Mathieu moonshine from both sides}
\label{sec:moonshine}
%%%%%%%%%%%%%%%%%%%%%%%%%%%%%%%%%%%%%%%%%%%%%%%%%%%%%%%%%%%%%%%%%%%%%%%%%%%%%%%

\subsection{The modular side}
Cheng--Harrison~\cite{ChengHarrison2014} (\texttt{1406\_0619.txt} ll.~800--813, eqs.~(3.5)--(3.6))
give $H^{(2)}(\tau)=(-2E_2(\tau)+48F_2^{(2)}(\tau))/\eta(\tau)^3$ with
$F_2^{(2)}=\sum_{r>s>0,\ r-s\text{ odd}}(-1)^r s\,q^{rs/2}$. \tierL.

\begin{theorem}[Computed equals typed] \tierA\
The series computed from this formula, divided exactly by $\prod(1-q^n)^3$, equals
$-2,2A_1,\dots,2A_9$ of the EOT table (\lean{hComputed\_eq\_table}); the transcription of $F_2$ is
checked separately (\lean{f2\_control}).
\end{theorem}

Cheng--Duncan--Harvey (CDH)~\cite{CDH2012} give the twined series $H_g=(\chi_g/24)H+F_g/\eta^3$
(\texttt{1204\_2779.txt} l.~2934, eq.~(4.18)) with $F_g$ the weight-2 forms of their Table~3
(ll.~2882--2915): multiples of $\Lambda_N$, eta quotients, and the newforms
$f_{11},f_{14},f_{15},f_{23,a},f_{23,b}$ with prefactors $2/5,1/3,1/4,1/3,1/11$. \tierL.

\begin{theorem}[All 21 columns of CDH Table~20] \tierA\
For every conjugacy class, the twined series computed as $24D\cdot H_g$ over $\mathbb{Z}$ equals $24D$
times the printed column of CDH Table~20 (ll.~4494--4512) through $q^9$: \lean{twined\_2A},
\lean{twined\_3A}, \lean{twined\_5A}, \lean{twined\_7AB} and \lean{twined\_2B}, \dots,
\lean{twined\_23AB}; the division by $24D$ is exact (\lean{twined\_div24}, \lean{twinedAll\_div}).
\end{theorem}
Negative controls: dropping $F_g$ breaks the match (\lean{twined\_2A\_needs\_F}); dropping the newforms
breaks $11A$ and $23AB$ (\lean{twined\_11A\_needs\_newform}, \lean{twined\_23AB\_needs\_newforms}); a
wrong exponent in the $12A$ eta quotient breaks it (\lean{twined\_12A\_exponent\_matters}). The two
forms that CDH print both as $\Lambda$-combinations and as eta quotients agree
(\lean{eta\_lambda\_agree\_2B}, \lean{eta\_lambda\_agree\_4A}).

\subsection{The group side}
CDH Table~8 (ll.~4170--4209) is the character table of $\Mtf$, with irrational values
$b_n=(-1+\sqrt{-n})/2$ for $n=7,15,23$. We store $a+b\,b_n$ as a pair $(a,b)$ with $b_n^2=-b_n-(n+1)/4$,
and accumulate inner products as a tuple (rational part, $\omega_7$-, $\omega_{15}$-, $\omega_{23}$-part);
since $1,\sqrt{-7},\sqrt{-15},\sqrt{-23}$ are $\mathbb{Q}$-linearly independent, an inner product equals
an integer exactly when the tuple is (integer, 0, 0, 0).

\begin{theorem}[Guards on the transcription] \tierA\
The column norms of the transcribed table are the centralizer orders (\lean{colNorm\_eq\_cent}), each
divisible by the element order and dividing $|\Mtf|$ (\lean{cent\_divides}); the class equation holds
(\lean{class\_equation}); the first orthogonality relation holds for all $26\times26$ pairs of rows,
irrational columns included (\lean{gram\_ok}); a misreading outside the relabelling orbit violates it
(\lean{misread\_fails\_gram}).
\end{theorem}
The flattened PDF text loses the overlines ($b_7$ and $\bar b_7$ both print as \texttt{b7}). An
exhaustive search (Python, \emph{not} a kernel proof) found 128 readings passing orthogonality, exactly
the orbit under renaming classes and characters, all with the same traces.

\begin{theorem}[Moonshine at every class] \tierA\
For all 26 classes $g$ and levels $n=1,\dots,7$, the trace of $g$ on EOT's representation is a rational
integer (the irrational parts vanish) and equals the coefficient of the computed $H_g$
(\lean{trace\_eq\_twined\_coeff\_all}); algebraically conjugate classes have equal traces
(\lean{trace\_pairs\_equal}); swapping two classes breaks the match (\lean{trace\_7A\_ne\_series\_23A}).
\end{theorem}

\subsection{Levels 0--9: the module, from computed data}
Reversing the direction, the multiplicity of each irreducible in the graded piece $K_n$ is
$\langle T_n,\chi_i\rangle$ with $T_n(g)$ the computed coefficient. CDH print the decompositions in their
Table~48 (ll.~5305--5330). \tierL.

\begin{theorem}[The first ten graded pieces] \tierA\
For $n=0,\dots,9$ every multiplicity is an integer, equal to CDH Table~48, and non-negative from $n=1$
(\lean{moonshine\_modules\_decompose}, \lean{moonshine\_multiplicities\_nonneg}). In particular levels
$8$ and $9$, for which EOT give no decomposition, are characters of genuine $\Mtf$-modules; e.g.
$K_8=990\oplus\overline{990}\oplus1035'\oplus1035''\oplus2\cdot1265\oplus\cdots\oplus6\cdot10395$.
\end{theorem}
Cheng~\cite{Cheng2010} (\texttt{1005\_5415.txt} ll.~540--548) observed that EOT's original proposal
$30843=10395+2\cdot5796+5544+3312$ at level~7 is inconsistent with the twined series and proposed
$10395+5796+5544+5313+2024+1771$; the first is refuted by \lean{eot\_level7\_proposal\_inconsistent},
the second is the decomposition used in \lean{trace\_eq\_twined\_coeff\_all}. \tierA.

%%%%%%%%%%%%%%%%%%%%%%%%%%%%%%%%%%%%%%%%%%%%%%%%%%%%%%%%%%%%%%%%%%%%%%%%%%%%%%%
\section{The forger's test}
\label{sec:forger}
%%%%%%%%%%%%%%%%%%%%%%%%%%%%%%%%%%%%%%%%%%%%%%%%%%%%%%%%%%%%%%%%%%%%%%%%%%%%%%%
Paper~7's Theorem~6.1 relates the first two coefficients, $\mathcal{A}_2/(N_Q\mathcal{A}_1)=462/360$,
with product $27720$. It holds at the identity, now from the computed series (\lean{lock\_at\_identity}).

\begin{theorem}[The lock fails twining] \tierA\
The ratio form of the lock, with each twined series' own coefficients, fails at all 25 non-identity
classes (\lean{ratio\_fails\_at\_every\_class}; first shown at $2A,3A,5A,7AB$ by
\lean{ratio\_fails\_at\_2A}, \dots). At $2A$, $\mathcal{A}_2/\mathcal{A}_1=14/(-6)$, not $462/90$.
\end{theorem}
Relations that encode module structure survive twining; this ratio does not. By that criterion the lock
is numerology. \tierC\ (the criterion is a reading; the failure is \tierA).

By contrast, the integers of the dyon sector pass. CDH's Frame shapes (Table~14, ll.~4349--4364) have
degree 24 and fixed points $\chi_g$ (\lean{frame\_degree}, \lean{frame\_fixed\_points}) and agree with the
power maps of Table~8 (\lean{frame\_power\_maps}); the twined Göttsche numbers
$T_k(g)=[p^k]\prod_{n\ge1}\det(1-p^ng\,|\,\mathbb{C}^{24})^{-1}$ equal $p_{24}(k)$ at $g=1$
(\lean{twined\_goettsche\_identity}) and, for $k\le4$, are characters of $\Mtf$ with non-negative
integer multiplicities (\lean{twined\_goettsche\_is\_character}); e.g.
$H^*(\mathrm{Hilb}^2K3)=3\cdot\mathbf{1}\oplus3\cdot\mathbf{23}\oplus\mathbf{252}$
(\lean{twined\_goettsche\_k2}). \tierA. This is mathematically expected (symmetric powers of a
permutation module); its value is the certified cross-check of three transcribed tables and the contrast
with the lock.

%%%%%%%%%%%%%%%%%%%%%%%%%%%%%%%%%%%%%%%%%%%%%%%%%%%%%%%%%%%%%%%%%%%%%%%%%%%%%%%
\section{Where the 24 of the shadow comes from}
\label{sec:shadow}
%%%%%%%%%%%%%%%%%%%%%%%%%%%%%%%%%%%%%%%%%%%%%%%%%%%%%%%%%%%%%%%%%%%%%%%%%%%%%%%
Harvey--Murthy--Nazaroglu (HMN)~\cite{HMN2014} state that the shadow of $H^{(2)}$ is $24\eta^3$
(\texttt{1410\_6174.txt} l.~246), and that $\eta^3H$ is a mixed mock modular form whose completion is,
up to $-\tfrac12$, the second helicity supertrace of the double-scaled little string theory at $k=2$
(ll.~262--270). \tierL. No completion or modular transformation is formalized here. What is formalized
is the arithmetic skeleton, following CDH §2.3 (\texttt{1204\_2779.txt} ll.~657--711, eqs.~(2.19)--(2.25)):
$\psi=\Psi_{1,1}\,Z=\chi\,\mathrm{Av}^{(2)}[(y+1)/(y-1)]+H\,\hat\theta_1$, with $Z=8(f_2^2+f_3^2+f_4^2)$
((2.49), (2.62), ll.~1266, 1397), cleared of denominators into an identity of $q$-series with exact
Laurent-polynomial coefficients in $y$.

\begin{theorem}[The polar/finite decomposition] \tierA\
Through $q^9$: $Z(\tau,0)=24$ with every higher coefficient zero (\lean{ellipticGenus\_z0}); $Z$ begins
$2y+20+2y^{-1}+q(20y^2-128y+216-128y^{-1}+20y^{-2})$ (\lean{ellipticGenus\_first\_terms}); the cleared
$\Psi_{1,1}$ reproduces CDH's printed expansion (\lean{psi11\_expansion}); the decomposition holds with
$\chi=24$ and $H$ the series computed from $(-2E_2+48F_2)/\eta^3$ (\lean{decomposition}), and fails for
$\chi=23,25$ (\lean{decomposition\_pins\_chi}) and for $H$ replaced by $H_{2A}$
(\lean{decomposition\_needs\_H}); the unary theta series whose multiple is the shadow equals $\eta^3$
(\lean{shadowTheta\_eq\_eta3}); CDH's twined shadow multiplicities are traces on $\mathbf{1}\oplus\mathbf{23}$
(\lean{shadow\_coeff\_eq\_perm\_trace}).
\end{theorem}
Two independent facts meet at $24$: the value $Z(\tau,0)$ and the multiplicity of the polar part. That
this multiplicity is the shadow's coefficient is Zwegers' theorem as recast by CDH. \tierL. Reading the
polar (massless) and finite (massive) parts as ``macro'' and ``micro'' sectors is \tierC.

%%%%%%%%%%%%%%%%%%%%%%%%%%%%%%%%%%%%%%%%%%%%%%%%%%%%%%%%%%%%%%%%%%%%%%%%%%%%%%%
\section{Double-scaled little strings}
\label{sec:hmn}
%%%%%%%%%%%%%%%%%%%%%%%%%%%%%%%%%%%%%%%%%%%%%%%%%%%%%%%%%%%%%%%%%%%%%%%%%%%%%%%
HMN attach to each ADE root system $Y$ with Coxeter number $k$ (the lattice of vanishing $(-2)$-cycles of
$k$ NS5-branes) a BPS index $\chi_2^Y=\mathrm{rk}(Y)E_2-24F_2^Y$ (\texttt{1410\_6174.txt} l.~732,
eq.~(2.36)), with $F_2^{(k,d)}$ an explicit divisor sum (l.~703, eq.~(2.29)) combined according to their
Table~1, and relate it to umbral moonshine for $X=Y^{24/\mathrm{rk}(Y)}$ by
$\chi_2^Y=-(\mathrm{rk}(Y)/2)\chi_2^X$ (l.~1266, eq.~(4.9)). \tierL.

\begin{theorem}[HMN's formula, reproduced] \tierA\
The formula reproduces HMN's printed series for $Y=A_3$ and $Y=A_7$ (ll.~1295, 1299;
\lean{hmn\_A3\_printed}, \lean{hmn\_A7\_printed}); HMN's $F_2^{(2,1)}$ equals Cheng--Harrison's $F_2$
through $q^{40}$ (\lean{f2kd\_k2\_eq\_f2Coeff}); and at $k=2$, $-2\chi_2^{A_1}=\eta^3H=-2E_2+48F_2$ through
$q^9$ (\lean{hmn\_k2\_eq\_moonshine}).
\end{theorem}

\begin{theorem}[The umbral relation (4.9)] \tierA\
For $\ell=2,3,4,5,7,13$ the umbral Jacobi forms $Z^{(\ell)}=2\phi_1^{(\ell)}$ built from theta functions
(CDH (2.49)--(2.51)) are integral (\lean{umbral\_exact}), satisfy $Z^{(\ell)}(\tau,0)=24/(\ell-1)$
(\lean{umbral\_chi}), their polar part removes the pole (\lean{umbral\_pole\_removed},
\lean{umbral\_pole\_removed\_13}), and HMN's (4.9) holds against $Y=A_{\ell-1}$ through $q^6$
(\lean{hmn\_umbral\_relation\_2}, \dots, \lean{hmn\_umbral\_relation\_13}); pairing $\ell=3$ with $A_3$
fails (\lean{hmn\_umbral\_relation\_wrong\_Y}).
\end{theorem}

HMN observe (ll.~1284--1299) that for umbral $Y$ all coefficients of $\chi_2^Y$ are divisible by
$\mathrm{rk}(Y)$, ``but it is not clear why'' for the massive states, and that $A_7$ fails.

\begin{theorem}[The divisibility, both directions] \tierA\
For every $Y$ and every $n$, $\mathrm{rk}(Y)\mid24$ implies $\mathrm{rk}(Y)\mid[q^n]\chi_2^Y$
(\lean{divisible\_of\_dvd\_24}; all coefficients, not a truncation, since
$[q^n]\chi_2^Y=-24(\mathrm{rk}(Y)\sigma(n)+F_2^Y(n))$ with $F_2^Y(n)\in\mathbb{Z}$). Conversely, for every
$A_{k-1}$ ($k\ge2$) and every $D_{j+1}$ ($j\ge3$), all coefficients are divisible by the rank \emph{iff}
the rank divides 24, with the $q^1$ coefficient as witness (\lean{A\_divisible\_iff},
\lean{D\_divisible\_iff}); $E_6,E_8$ divide and $E_7$ fails (\lean{E\_divisibility}).
\end{theorem}
The divisible $Y$ are exactly HMN's umbral list. Within HMN's own formula, the puzzle has an elementary
cause: the massive coefficients are $24\times$ integers. Why the factor in (2.36) is $24$ is not addressed;
reading it as $\chi(K3)$ is \tierC.

\begin{casebox}{Closed as not available: the link to $(-2)$-reflections}
The Stream~4 plan asked to relate this counting to the kernel-checked fact that reflections in
$(-2)$-vectors are lattice isometries (paper~8), or to show that no formal relation is available. In HMN
the lattice-theoretic structure enters through the Niemeier lattice $L_X$ and its Weyl group,
$G^X=\mathrm{Aut}(L_X)/W_X$ (\texttt{1410\_6174.txt} l.~277, eq.~(1.10)). Nothing computed above involves a
reflection or an isometry; $\mathrm{rk}(Y)$ enters only as a number. Constructing $L_{A_1^{24}}$ would not
close the gap, which is the isomorphism $\mathrm{Aut}(L_X)/W_X\cong\Mtf$; we record the link as not
available rather than implying a closure.
\end{casebox}

%%%%%%%%%%%%%%%%%%%%%%%%%%%%%%%%%%%%%%%%%%%%%%%%%%%%%%%%%%%%%%%%%%%%%%%%%%%%%%%
\section{\texorpdfstring{Dyons on $K3\times T^2$}{Dyons on K3 x T2}}
\label{sec:dyons}
%%%%%%%%%%%%%%%%%%%%%%%%%%%%%%%%%%%%%%%%%%%%%%%%%%%%%%%%%%%%%%%%%%%%%%%%%%%%%%%

\subsection{The partition function from the elliptic genus}
Dabholkar--Murthy--Zagier (DMZ)~\cite{DMZ2012}: quarter-BPS dyons of type~II string theory on
$K3\times T^2$ are counted by $1/\Phi_{10}(\Omega)=\sum_{m\ge-1}\psi_m(\tau,z)p^m$ (\texttt{1208\_4074.txt}
l.~177, eq.~(1.1)); $\Phi_{10}$ is the Borcherds (multiplicative) lift of the $K3$ elliptic genus
$2\phi_{0,1}=\sum c(4n-l^2)q^ny^l$ (l.~1820, eq.~(5.13)). Hence $\Delta\psi_m=A^{-1}G_{m+1}$ with
$A=\phi_{-2,1}$ and $\sum_kG_kp^k=\prod_{r\ge1,s\ge0,t}(1-p^rq^sy^t)^{-c(4rs-t^2)}$. \tierL.

\begin{theorem}[The product, from our $Z_{K3}$] \tierA\
The coefficients of the computed $Z_{K3}$ depend only on $4n-l^2$ (\lean{c\_depends\_only\_on\_D});
every $c(D)$ the truncated products read is available through $q^9$ (\lean{dmvv\_reachable}, an
enumerated guard, with a failing example \lean{dmvv\_unreachable\_example}); the integer divisions in the
product are exact (\lean{negBinom\_exact}). At $z=0$, $G_k=p_{24}(k)=1,24,324,3200,25650,176256$ for
$k\le5$, constant in $\tau$ (\lean{goettsche}, \lean{p24\_values}) --- the Euler numbers of
$\mathrm{Hilb}^kK3$, with $p_{24}$ the 24-coloured partitions (DMZ l.~1070; $p_{24}(3)=3200$ at l.~3639).
\end{theorem}

\begin{theorem}[DMZ (5.16)] \tierA\
The six printed identities $\Delta\psi_{-1}=A^{-1}$, \dots,
$72\Delta\psi_4=51A^{-1}B^5+155E_4AB^3+93E_6A^2B^2+102E_4^2A^3B+31E_4E_6A^4$ (\texttt{1208\_4074.txt}
ll.~1879--1888), with $A$, $B=\phi_{0,1}$ from theta products and $E_4,E_6$ from divisor sums, hold
(\lean{dmz\_516\_q1}: $k\le5$ through $q^1$; \lean{dmz\_516\_q2}: $k\le4$ through $q^2$).
\end{theorem}

\subsection{Two-centred and single-centred parts}
DMZ's canonical decomposition $\psi_m=\psi_m^F+\psi_m^P$ has polar part
$\psi_m^P=p_{24}(m+1)\Delta^{-1}A_{2,m}$, $A_{2,m}=\sum_sq^{ms^2+s}y^{2ms+1}/(1-q^sy)^2$ (ll.~196--204,
eqs.~(1.2)--(1.3)), counting two-centred configurations that appear and disappear across walls of
marginal stability, while $\psi_m^F$ counts single-centred (``immortal'') black holes (ll.~424--443).
\tierL. $A_{2,m}$ has different Fourier expansions in different strips; all statements below use
$|q|<|y|<1$.

\begin{theorem}[The polar part] \tierA\
In the strip, the direct expansion of (1.3) equals DMZ's (9.55) (l.~5210; \lean{a2m\_strip\_eq\_955});
$G_{m+1}-p_{24}(m+1)A\,A_{2,m}$ vanishes to second order at $y=1$ at every order in $q$, $m=1,2,3$
(\lean{polar\_part\_removes\_pole}), and $p_{24}(m+1)\pm1$ do not (\lean{polar\_coefficient\_pinned}).
\end{theorem}
The second-order condition is automatic by $y\leftrightarrow y^{-1}$ symmetry; what these two theorems
discriminate is the first-order condition --- the pole coefficient is $G_{m+1}(\tau,0)=p_{24}(m+1)$. The
full second-order content (exact division by $A$) is carried by the theorems below.

\begin{theorem}[Immortal dyons are class numbers] \tierA\
With $H$ the generating function of Hurwitz class numbers, computed by counting reduced binary quadratic
forms (\lean{hurwitz\_values}), the single-centred counting functions computed from the product satisfy,
through the orders stated,
\[
\Delta\psi_1^F=3E_4A-648H\ \ (q^3),\qquad
3\Delta\psi_2^F=-800\,(12H|V_2)+22E_4AB-10E_6A^2\ \ (q^2),
\]
\[
48\Delta\psi_3^F=-102600\,(12H|V_3)+467E_4AB^2-430E_6A^2B+203E_4^2A^3\ \ (q^2),
\]
where $12(H|V_p)$ has coefficients $12H(\Delta)+p\cdot12H(\Delta/p^2)$
(\lean{immortal\_m1}, \lean{immortal\_m2}, \lean{immortal\_m3}; exactness \lean{immortal\_m1\_exact},
\lean{immortal\_exact\_m23}). Dropping $H$ breaks $m=1$ (\lean{immortal\_m1\_needs\_H}); dropping the
Hecke correction breaks $m=2$ (\lean{immortal\_m2\_needs\_V2}).
\end{theorem}
The combinations follow from (5.16) and DMZ's optimal forms: Example~5 (ll.~3530--3558), (9.10) (l.~3774),
(9.11) (l.~3793), (9.12)--(9.13) (ll.~3830, 3850). These are checked, not cited: the finite parts of
$\psi_{0,3}^{\mathrm{opt}}/A$ and $\psi_{0,4}^{\mathrm{opt}}/A$ are $-12^{m+1}H|V_m$
(\lean{dmz\_911\_verified}, \lean{dmz\_913\_verified\_m3}), and DMZ's printed tables are reproduced
(\lean{example5\_tables}, \lean{dmz\_911\_table}, \lean{dmz\_912\_table}). \tierA.

%%%%%%%%%%%%%%%%%%%%%%%%%%%%%%%%%%%%%%%%%%%%%%%%%%%%%%%%%%%%%%%%%%%%%%%%%%%%%%%
\section{\texorpdfstring{$\Mtf$-twisted dyons}{M24-twisted dyons}}
\label{sec:twisted}
%%%%%%%%%%%%%%%%%%%%%%%%%%%%%%%%%%%%%%%%%%%%%%%%%%%%%%%%%%%%%%%%%%%%%%%%%%%%%%%
Cheng~\cite{Cheng2010} proposes twisted denominators
$\Phi_g=e^{4\pi i(\rho,\Omega)}\prod_\alpha\exp\big(-\sum_{k\ge1}\tfrac1kc_{g^k}(|\alpha|^2)e^{2\pi ik(\alpha,\Omega)}\big)$
built from the twisted elliptic genera $Z_{g^k}$ (\texttt{1005\_5415.txt} ll.~933--942, eqs.~(3.9)--(3.10)),
with a pole factorizing into two twisted half-BPS functions $1/\eta_g$ (l.~985, eq.~(3.11)); for the
geometric classes $Z_{g_p}=\tfrac2{p+1}\phi_{0,1}+\dots$ (l.~423, eq.~(2.5)). \tierL. We compute
$Z_g=(\chi_g/24)Z-F_g\,A$ for all 26 classes from Stream~4's data (a consequence of CDH (4.18) and the
decomposition of \Cref{sec:shadow}), and the product by Newton's identities with CDH's power maps.

\begin{theorem}[Twisted dyons are virtual characters] \tierA\
The 26 twisted genera are integral, index~1 and satisfy $Z_g(\tau,0)=\chi_g$ (\lean{twined\_genus\_exact},
\lean{twined\_genus\_index1}, \lean{twined\_genus\_z0}); Newton's divisions are exact
(\lean{newton\_exact}); $g=1A$ recovers the untwisted product (\lean{twisted\_product\_untwisted}); at $z=0$
the $p$-coefficients are the twined Göttsche numbers (\lean{twisted\_goettsche}); every coefficient of
$G_1^{(g)},G_2^{(g)}$ through $q^2$ is a virtual character of $\Mtf$ (\lean{twisted\_dyons\_are\_characters});
for all 26 classes $T_2(g)\,A_{2,1}$ removes the double pole (\lean{twisted\_polar\_removes\_pole}); and the
twisted single-centred counting function at $m=1$ has virtual-character coefficients
(\lean{twisted\_immortal\_are\_characters}), e.g. the $q^1y^0$ coefficient
$-1800=2\cdot\mathbf1+2\cdot\mathbf{23}+2\cdot\mathbf{45}+2\cdot\overline{\mathbf{45}}-\mathbf{231}
-\overline{\mathbf{231}}+2\cdot\mathbf{252}-\mathbf{1035}'-\mathbf{1035}''$ (\lean{twisted\_immortal\_example}).
\end{theorem}
Once the twisted genera are characters (Gannon's theorem, \tierL), the later steps are linear and
class-independent, so character-valuedness is expected; the content is the certified end-to-end
computation from modular data and CDH's power maps. That $1/\Phi_g$ counts twisted dyons, and any
$\Mtf$ action on black-hole microstates, is \tierL/\tierC.

%%%%%%%%%%%%%%%%%%%%%%%%%%%%%%%%%%%%%%%%%%%%%%%%%%%%%%%%%%%%%%%%%%%%%%%%%%%%%%%
\section{Reading notes on the literature}
\label{sec:notes}
%%%%%%%%%%%%%%%%%%%%%%%%%%%%%%%%%%%%%%%%%%%%%%%%%%%%%%%%%%%%%%%%%%%%%%%%%%%%%%%
These are recorded as observations, not corrections; each is backed by a kernel-checked computation.
\begin{enumerate}
\item \textbf{HMN (4.2).} \texttt{1410\_6174.txt} l.~1194 prints $\chi_2^Y=-\mathrm{rk}(Y)E_2+24F_2^Y$,
the opposite overall sign of their (2.36) (l.~732). Their printed series (4.10)--(4.11) (ll.~1295, 1299)
and the count of $\mathrm{rk}(Y)$ massless vector multiplets require the sign of (2.36), which is the one
used (\lean{hmn\_A3\_printed}, \lean{hmn\_A7\_printed}).
\item \textbf{DMZ (1.4).} \texttt{1208\_4074.txt} l.~211 lists, for the strip $0<\mathrm{Im}\,z<\mathrm{Im}\,\tau$,
the terms $\ell q^{(r^2-\ell^2)/4m}y^r$ with $r\ge\ell>0$. The direct expansion of (1.3) in that strip, and
DMZ's own general formula (9.55), also contain the mirror terms with $y^{-r}$, $r>\ell>0$ (from $s\le-1$);
the formalization uses (9.55) (\lean{a2m\_strip\_eq\_955}).
\item \textbf{Cheng's $c_g(-1)$.} \texttt{1005\_5415.txt} l.~960 writes $c_g(-1)=-2$. In Cheng's own
(2.5) and in the normalization used here the coefficient of $y^{\pm1}q^0$ is $+2$ for all 26 classes
(\lean{twisted\_genus\_q0}, \lean{c\_minus\_one}); the printed $-2$ is a sign convention or a misprint.
\item \textbf{EOT's level-7 decomposition.} Already noted by Cheng (ll.~540--548); confirmed by
\lean{eot\_level7\_proposal\_inconsistent}.
\end{enumerate}

%%%%%%%%%%%%%%%%%%%%%%%%%%%%%%%%%%%%%%%%%%%%%%%%%%%%%%%%%%%%%%%%%%%%%%%%%%%%%%%
\section{Context: the programme's experimental verdict}
\label{sec:context}
%%%%%%%%%%%%%%%%%%%%%%%%%%%%%%%%%%%%%%%%%%%%%%%%%%%%%%%%%%%%%%%%%%%%%%%%%%%%%%%
Streams~4 and~5 study the mathematics of $\mathcal{N}=4$ string compactifications on $K3\times T^2$; they
make no prediction about our universe. The programme's single experimental prediction came from
Stream~3 (paper~9~\cite{Callens2026Paper9}): the self-dual length $s=\sqrt{\ell_P\,c/H_0}\approx47\,\mu$m
read as the radius of one large extra dimension. In Stream~6 that prediction was frozen (git tag
\texttt{stream6-p1-frozen}; the freeze was made after the bounds were known, and is disclosed as a
retrodiction) and compared with the Eöt-Wash 2020 bounds~\cite{LeeAdelberger2020} (toroidal radius
$<30\,\mu$m, Yukawa range $<38.6\,\mu$m; \texttt{2002\_11761.txt} ll.~285--289) and the neutron-star bound
of Montero--Vafa--Valenzuela~\cite{MonteroVafaValenzuela2023} ($<44\,\mu$m; \texttt{2205\_12293.txt}
l.~327): it fails all three and is excluded by the pre-registered rule
(\lean{DualScaleCosmology.Stream6Verdict.p1\_verdict\_excluded}); the programme's own T-duality fixes the
order-one factor at $1$ (\lean{p62\_selfdual\_fixed\_point}), so no convention rescues it. The results of
this paper are unaffected: they are Tier~A mathematics, and the ``micro/macro'' readings above were
labelled Tier~C throughout.

%%%%%%%%%%%%%%%%%%%%%%%%%%%%%%%%%%%%%%%%%%%%%%%%%%%%%%%%%%%%%%%%%%%%%%%%%%%%%%%
\section{What is, and is not, proved}
\label{sec:scope}
%%%%%%%%%%%%%%%%%%%%%%%%%%%%%%%%%%%%%%%%%%%%%%%%%%%%%%%%%%%%%%%%%%%%%%%%%%%%%%%
\begin{scopebox}
\textbf{Is proved (Tier~A; \texttt{DualScaleMoonshine} 101 theorems, \texttt{DualScaleDyons} 44 theorems,
0 failing, 0 \texttt{sorry}; axioms at most \texttt{propext}, \texttt{Classical.choice}, \texttt{Quot.sound}):}
finite identities between truncated series, each through the order stated in its theorem ---
$H^{(2)}$ and all 21 twined series from their formulas; consistency guards and full orthogonality of the
transcribed character table; traces equal computed coefficients at all 26 classes (levels 1--7);
multiplicities of $K_0,\dots,K_9$ equal CDH Table~48; the lock's failure at 25 classes; the polar/finite
decomposition of $\Psi_{1,1}Z$ with multiplicity $24$; HMN's formula and umbral relation for six
lambencies; the divisibility theorem in both directions for all $A$ and $D$ types; the dyon product from
$Z_{K3}$, DMZ's (5.16), the polar coefficient $p_{24}(m+1)$, the immortal parts at $m=1,2,3$ via class
numbers; the twisted dyon products and their character-valued coefficients.

\textbf{Is not proved:} mock modularity of any series, any shadow, completion or modular transformation;
anything past the stated truncations; the Siegel-modularity of $\Phi_{10}$ or convergence of the product;
that $1/\Phi_{10}$ or $1/\Phi_g$ counts dyons, or that $\psi_m^F$ counts single-centred black holes; the
existence of the moonshine module (Gannon); the general statements behind the finite checks (e.g. DMZ's
(9.13) for all primes); that the transcribed character table \emph{is} the character table of $\Mtf$ (the
guards are consistency checks it must pass, not a construction of $\Mtf$); the link to
$(-2)$-reflections (\Cref{sec:hmn}); and any statement about our universe.
\end{scopebox}

%%%%%%%%%%%%%%%%%%%%%%%%%%%%%%%%%%%%%%%%%%%%%%%%%%%%%%%%%%%%%%%%%%%%%%%%%%%%%%%
\section{Methods}
\label{sec:methods}
%%%%%%%%%%%%%%%%%%%%%%%%%%%%%%%%%%%%%%%%%%%%%%%%%%%%%%%%%%%%%%%%%%%%%%%%%%%%%%%
\textbf{Exact arithmetic.} Series are lists of integers; Jacobi-form coefficients are Laurent polynomials
stored as (offset, coefficient list), with exact convolution --- no truncation in $y$. Denominators are
cleared by hand and the cleared identities are what the kernel checks; exact divisions are asserted by
separate theorems. Every formula was prototyped numerically before being stated in Lean.

\textbf{Kernel evaluation.} All decisions use \texttt{decide}; the heaviest (umbral forms at
$\ell=13$, twisted dyons) use \texttt{decide +kernel}, in which the kernel evaluates the decision procedure
directly --- no compiled evaluation and no additional axiom. The HMN file takes about 17 minutes to check,
the twisted-dyon file about 8.

\textbf{Guards against silent truncation.} A truncated product silently reads unavailable coefficients as
zero. The reachability of every coefficient read by the dyon product is an enumerated theorem with a
failing example, and every negative control in this paper is a theorem that the corresponding identity is
\emph{false}.

\textbf{Two errors caught by computation.} An off-by-one in the indexing of $p_{24}$ (DMZ use the number of
24-coloured partitions of $n$, so $p_{24}(2)=324$) made a first prototype of the polar-part check fail for
every coefficient, including the correct one; the failure of the positive case is what exposed it. A
reachability estimate stated as ``$KQ\le9$'' was replaced by an enumeration after review, because the
estimate was not connected to the computation it guarded.

%%%%%%%%%%%%%%%%%%%%%%%%%%%%%%%%%%%%%%%%%%%%%%%%%%%%%%%%%%%%%%%%%%%%%%%%%%%%%%%
\section{Open problems for the community}
\label{sec:open}
%%%%%%%%%%%%%%%%%%%%%%%%%%%%%%%%%%%%%%%%%%%%%%%%%%%%%%%%%%%%%%%%%%%%%%%%%%%%%%%
\begin{enumerate}
\item \textbf{Mock modularity in a proof assistant.} No statement here involves a completion or a modular
transformation. Formalizing Zwegers' completion of an Appell--Lerch sum (even for one example) would turn
the Tier~L step ``polar multiplicity $=$ shadow coefficient'' into Tier~A.
\item \textbf{The general theorems behind the finite checks.} DMZ's (9.13) ($\Phi^{\mathrm{opt}}_{2,p}=-H|V_p$
for all primes $p$) and HMN's (4.9) for all umbral $Y$ are checked here at small order and small index only.
\item \textbf{$D$- and $E$-type umbral forms.} The umbral relation is checked for pure $A$-type Niemeier root
systems; the $D$ and $E$ cases need umbral forms not constructed from the theta products used here.
\item \textbf{The reflection bridge.} Formalizing $\mathrm{Aut}(L_X)/W_X\cong\Mtf$ for $X=A_1^{24}$, or any
Niemeier lattice, would reopen the link closed as not available in \Cref{sec:hmn}.
\item \textbf{Twisted immortal counts at higher order.} The character-valuedness of the twisted
single-centred counts is checked through $q^2$ at $m=1$; kernel cost grows quickly with $m$ and $q$.
\item \textbf{Why 24 in HMN's (2.36)?} The divisibility theorem reduces HMN's puzzle to the factor 24; a
derivation of that factor from the DSLST on $K3$ would complete the explanation.
\end{enumerate}

%%%%%%%%%%%%%%%%%%%%%%%%%%%%%%%%%%%%%%%%%%%%%%%%%%%%%%%%%%%%%%%%%%%%%%%%%%%%%%%
\section{Proposed research directions}
\label{sec:directions}
%%%%%%%%%%%%%%%%%%%%%%%%%%%%%%%%%%%%%%%%%%%%%%%%%%%%%%%%%%%%%%%%%%%%%%%%%%%%%%%
\begin{enumerate}
\item \textbf{A library of computed moonshine.} The same method (closed formula, twined formula, character
table, kernel comparison) applies to the other umbral groups; CDH's tables for $\ell=3,4,5,7,13$ are pinned
and would serve as transcription controls.
\item \textbf{The forger's test as a standard.} Any integer coincidence proposed as physics should first be
twined; formalizing the test for a proposed relation costs little once the twined series exist.
\item \textbf{Wall-crossing formalized.} The strip dependence of $A_{2,m}$ (DMZ (9.55)) is the arithmetic of
wall-crossing; checking the jump across a wall against the two-centred degeneracy formula would connect
the polar part to its physical interpretation at the level of identities.
\item \textbf{CHL models.} For the geometric classes, Cheng's twisted denominators coincide with CHL dyon
counting functions; checking this identification for $p=2,3,5,7$ would link the twisted products to an
independent physical computation.
\end{enumerate}

%%%%%%%%%%%%%%%%%%%%%%%%%%%%%%%%%%%%%%%%%%%%%%%%%%%%%%%%%%%%%%%%%%%%%%%%%%%%%%%
\section{Conclusion}
\label{sec:conclusion}
%%%%%%%%%%%%%%%%%%%%%%%%%%%%%%%%%%%%%%%%%%%%%%%%%%%%%%%%%%%%%%%%%%%%%%%%%%%%%%%
Starting from the $K3$ elliptic genus computed from theta functions, and from printed tables used only as
controls, the kernel confirms Mathieu moonshine at every class and at the first ten levels, locates the
$24$ of the shadow, reproduces the arithmetic of double-scaled little strings with a two-sided divisibility
theorem, and builds the dyon partition function on $K3\times T^2$ whose single-centred part at $m=1,2,3$
is given by class numbers and whose twisted versions are character-valued. The number $24=\chi(K3)$ passes
every twining test applied to it; a coincidence proposed earlier in this programme does not. These are
results about mathematics, recorded with their scope; the programme's attempt to extract a prediction
about our universe from the same structure has been tested separately and found excluded.

%%%%%%%%%%%%%%%%%%%%%%%%%%%%%%%%%%%%%%%%%%%%%%%%%%%%%%%%%%%%%%%%%%%%%%%%%%%%%%%
\appendix
\section{Reproducibility}
\label{app:repro}
%%%%%%%%%%%%%%%%%%%%%%%%%%%%%%%%%%%%%%%%%%%%%%%%%%%%%%%%%%%%%%%%%%%%%%%%%%%%%%%

\subsection{Build}
\begin{center}
\footnotesize\ttfamily
lake build StringTheoryFoundation DualScaleM24Formalization DoubleFieldTheory \textbackslash\\
  DualScaleValidation Lean5Corpus StringTheoryFormalization DualScaleStream2 \textbackslash\\
  DualScaleCosmology DualScaleMoonshine DualScaleDyons
\end{center}
All ten libraries build without error (3799 jobs at the Stream~5 release; release tag \texttt{v3.14.0}).

\subsection{Axiom audit}
\begin{center}
\footnotesize\ttfamily
python3 tools/axiom\_audit.py DualScaleMoonshine\\
python3 tools/axiom\_audit.py DualScaleDyons
\end{center}
report \textbf{101} and \textbf{44 theorems, 0 failing}. The computational theorems depend on
\texttt{propext} only (six, across the two libraries, on no axiom at all); the divisibility proofs of \Cref{sec:hmn} use the three standard
axioms. Across the ten libraries of the repository: \textbf{611 theorems, 0 failing}.

\subsection{Statement lock}
\begin{center}
\footnotesize\ttfamily
python3 tools/statement\_lock.py --check \$(find DualScaleMoonshine DualScaleDyons -name '*.lean')
\end{center}
locks the declarations of both libraries (Stream~4: 9 files; Stream~5: 6 files); every addition was
checked to add entries only.

\subsection{Sources pinned}
Hash-pinned in \texttt{papers/foundations/MANIFEST.md}, line ranges quoted at the point of use:
{\small\texttt{1004.0956, 1406.0619, 1204.2779, 1410.6174, 1208.4074, 1005.5415, 2002.11761,
2205.12293}}.

\subsection{Plans}
\texttt{docs/STREAM4\_WORKFLOW.md} and \texttt{docs/STREAM5\_WORKFLOW.md} record every phase as closed
(\checkmark) or closed as not available with the reason; \texttt{docs/STREAM6\_EXPERIMENT\_PLAN.md} records
the experimental verdict.

\begin{thebibliography}{99}
\bibitem{Callens2026Paper7} X.~Callens and the SocrateAI Agora Collaboration, ``The Dual-Scale String Theory: Mechanized Foundations, Singularity Resolution, Mathieu Moonshine, and a Zero-Free-Parameter Cosmological Conjecture on $K3\times T^2$,'' \textit{SocrateAI Laboratory for Mathematical Physics \& Formal Verification} (2026).
\bibitem{Callens2026Paper9} X.~Callens and the SocrateAI Agora Collaboration, ``Testing a Micro/Macro Dual-Scale Hypothesis on $K3\times T^2$: Scale-Factor Duality, the Cohen--Kaplan--Nelson Bound, and the Dark-Energy Length in Lean~4,'' \textit{SocrateAI Laboratory for Mathematical Physics \& Formal Verification} (2026).
\bibitem{EOT2010} T.~Eguchi, H.~Ooguri, and Y.~Tachikawa, ``Notes on the K3 Surface and the Mathieu group $M_{24}$,'' \textit{Exper. Math.} 20 (2011) 91--96, \textit{arXiv:1004.0956}.
\bibitem{ChengHarrison2014} M.~C.~N.~Cheng and S.~Harrison, ``Umbral Moonshine and K3 Surfaces,'' \textit{Commun. Math. Phys.} 339 (2015) 221--261, \textit{arXiv:1406.0619}.
\bibitem{CDH2012} M.~C.~N.~Cheng, J.~F.~R.~Duncan, and J.~A.~Harvey, ``Umbral Moonshine,'' \textit{Commun. Number Theory Phys.} 8 (2014) 101--242, \textit{arXiv:1204.2779}.
\bibitem{HMN2014} J.~A.~Harvey, S.~Murthy, and C.~Nazaroglu, ``ADE Double Scaled Little String Theories, Mock Modular Forms and Umbral Moonshine,'' \textit{JHEP} 05 (2015) 126, \textit{arXiv:1410.6174}.
\bibitem{DMZ2012} A.~Dabholkar, S.~Murthy, and D.~Zagier, ``Quantum Black Holes, Wall Crossing, and Mock Modular Forms,'' \textit{arXiv:1208.4074} (2012).
\bibitem{Cheng2010} M.~C.~N.~Cheng, ``K3 Surfaces, N=4 Dyons, and the Mathieu Group M24,'' \textit{Commun. Num. Theor. Phys.} 4 (2010) 623--657, \textit{arXiv:1005.5415}.
\bibitem{LeeAdelberger2020} J.~G.~Lee, E.~G.~Adelberger, T.~S.~Cook, S.~M.~Fleischer, and B.~R.~Heckel, ``New Test of the Gravitational $1/r^2$ Law at Separations down to $52\,\mu$m,'' \textit{Phys. Rev. Lett.} 124 (2020) 101101, \textit{arXiv:2002.11761}.
\bibitem{MonteroVafaValenzuela2023} M.~Montero, C.~Vafa, and I.~Valenzuela, ``The Dark Dimension and the Swampland,'' \textit{JHEP} 02 (2023) 022, \textit{arXiv:2205.12293}.
\end{thebibliography}

\end{document}
